\chapter{2012年北京卷（文）}

\section{单选题}

\begin{problem}
已知集合 \(A = \{x \in \R \mid 3x + 2 > 0\}\)，\(B = \{x \in \R \mid (x + 1)(x - 3) > 0\}\)，则 \(A \cap B =\)（\quad）

\choices
    {\((-\infty, -1)\)}
    {\(\left(-1, -\frac{2}{3}\right)\)}
    {\(\left(-\frac{2}{3}, 3\right)\)}
    {\((3, +\infty)\)}
\end{problem}

\begin{problem}
在复平面内，复数 \(\frac{10}{3 + \i}\) 对应的点的坐标为（\quad）

\choices
    {(1,3)}
    {(3,1)}
    {\((-1,3)\)}
    {\((3, - 1)\)}
\end{problem}

\begin{problem}
设不等式组 \(\left\{ \begin{array}{l} 0 \leqslant x \leqslant 2, \\ 0 \leqslant y \leqslant 2 \end{array} \right.\) 表示的平面区域为 \(D\). 在区域 \(D\) 内随机取一个点，则此点到坐标原点的距离大于 2 的概率是（\quad）

\choices
    {\(\frac{1}{4}\)}
    {\(\frac{\pi - 2}{2}\)}
    {\(\frac{\pi}{6}\)}
    {\(\frac{4 - \pi}{4}\)}
\end{problem}

\begin{problem}
执行如图所示的程序框图，输出的 \(S\) 值为（\quad）

\choices
    {2}
    {4}
    {8}
    {16}
\end{problem}

\begin{problem}
函数 \( f(x) = x^{\frac{1}{2}} - \left(\frac{1}{2}\right)^x \) 的零点个数为（\quad）

\choices
    {0}
    {1}
    {2}
    {3}
\end{problem}

\begin{problem}
已知 \(\{a_{n}\}\) 为等比数列，下面结论中正确的是（\quad）

\choices
    {\(a_{1} + a_{3} > 2a_{2}\)}
    {\(a_1^2 + a_3^2 \geqslant 2a_2^2\)}
    {若 \(a_1 = a_3\)，则 \(a_1 = a_2\)}
    {若 \(a_{3} > a_{1}\)，则 \(a_{4} > a_{2}\)}
\end{problem}

\begin{problem}
某三棱锥的三视图如图所示，该三棱锥的表面积是（\quad）

\bitmapfigure[max width=0.76\linewidth,max totalheight=6.2cm]{img/2012/beijing_liberal/q7_fig1.png}

\choices
    {\(28 + 6\sqrt{5}\)}
    {\(30 + 6\sqrt{5}\)}
    {\(56 + 12\sqrt{5}\)}
    {\(60 + 12\sqrt{5}\)}
\end{problem}

\begin{problem}
某棵果树前 \(n\) 年的总产量 \(S_{n}\) 与 \(n\) 之间的关系如图所示. 从目前记录的结果看，前 \(n\) 年的年平均产量最高，\(m\) 的值为（\quad）

\choices
    {5}
    {7}
    {9}
    {11}
\end{problem}

\section{填空题}

\begin{problem}
直线 \( y = x \) 被圆 \( x^{2} + (y - 2)^{2} = 4 \) 截得的弦长为 \fillinblank{} \fillinblank{}.
\end{problem}

\begin{problem}
已知 \( \{a_{n}\} \) 为等差数列， \( S_{n} \) 为其前 n 项和，若 \( a_{1}=\frac{1}{2} \) ， \( S_{2}=a_{3} \) ，则 \( a_{2}= \) ； \( S_{n}= \)  \fillinblank{}.
\end{problem}

\begin{problem}
在 \(\triangle ABC\) 中，若 \(a = 3, b = \sqrt{3}, \angle A = \frac{\pi}{3}\)，则 \(\angle C\) 的大小为 \fillinblank{} \fillinblank{}.
\end{problem}

\begin{problem}
已知函数 \( f(x) = \lg x \). 若 \( f(ab) = 1 \)，则 \( f(a^2) + f(b^2) = \)  \fillinblank{}.
\end{problem}

\begin{problem}
已知正方形 \(ABCD\) 的边长为 1，点 \(E\) 是 \(AB\) 边上的动点，则 \(\overrightarrow{DE} \cdot \overrightarrow{CB}\) 的值为 \fillinblank{}；\(\overrightarrow{DE} \cdot \overrightarrow{DC}\) 的最大值为 \fillinblank{} \fillinblank{}.
\end{problem}

\begin{problem}
已知 \( f(x) = m(x - 2m)(x + m + 3), g(x) = 2^x - 2 \). 若 \( \forall x \in \R, f(x) < 0 \) 或 \( g(x) < 0 \)，则 \( m \) 的取值范围是 \fillinblank{} \fillinblank{}.
\end{problem}

\section{解答题}

\begin{problem}
已知函数 \( f(x) = \frac{(\sin x - \cos x)\sin 2x}{-\sin x} \). (1) 求 \( f(x) \) 的定义域及最小正周期； (2) 求 \( f(x) \) 的单调递减区间.
\end{problem}

\begin{problem}
如图 1，在 Rt\(\triangle\)ABC 中，\(\angle\)C = 90°，D, E 分别是 AC, AB 的中点，点 F 为线段 CD 上的一点，将 \(\triangle\)ADE 沿 DE 折起到 \(\triangle\)A\(_1\)DE 的位置，使 A\(_1\)F \(\perp\) CD，如图 2. (1) 求证: \( DE // \) 平面 \( A_{1}CB \) ; (2) 求证: \( A_{1}F \perp BE; \) (3) 线段 \(A_{1}B\) 上是否存在点 \(Q\)，使 \(A_{1}C \perp\) 平面 \(DEQ\)\(\perp\) 说明理由.

图1 图2
\end{problem}

\begin{problem}
近年来，某市为促进生活垃圾的分类处理. 将生活垃圾分为厨余垃圾、可回收物和其他垃圾三类，并分别设置了相应的垃圾箱. 为调查居民生活垃圾分类投放情况，现随机抽取了该市三类垃圾箱中总计 1000 吨生活垃圾，数据统计如下 (单位: 吨): “厨余垃圾”箱 “可回收物”箱 “其他垃圾”箱 厨余垃圾 400 100 100 可回收物 30 240 30 其他垃圾 20 20 60 (1) 试估计厨余垃圾投放正确的概率; (2) 试估计生活垃圾投放错误的概率; (3) 假设厨余垃圾在“厨余垃圾”箱，“可回收物”箱，“其他垃圾”箱的投放量分别为 a, b, c，其中 a > 0, \( a + b + c = 600 \) . 当数据 a, b, c 的方差 \( s^{2} \) 最大时，写出 a, b, c 的值（结论不要求证明），并求此时 \( s^{2} \) 的值. (注： \( s^{2}=\frac{1}{n}\left[(x_{1}-\overline{x})^{2}+(x_{2}-\overline{x})^{2}+\cdots+(x_{n}-\overline{x})^{2}\right] \) ，其中 \( \overline{x} \) 为数据 \( x_{1},x_{2},\cdots,x_{n} \) 的平均数)
\end{problem}

\begin{problem}
已知椭圆 \(C: \frac{x^2}{a^2} + \frac{y^2}{b^2} = 1 (a > b > 0)\) 的一个顶点为 \(A(2,0)\). 离心率为 \(\frac{\sqrt{2}}{2}\)，直线 \(y = k(x - 1)\) 与椭圆 \(C\) 交于不同的两点 \(M, N\). (i) 求椭圆 C 的方程; (2) 当 \(\triangle AMN\) 的面积为 \(\frac{\sqrt{10}}{3}\) 时，求 \(k\) 的值.
\end{problem}

\begin{problem}
设 \(A\) 是如下形式的 2 行 3 列的数表. a b c d e f 满足性质 \(P:a,b,c,d,e,f\in [-1,1]\) ，且 \(a + b + c + d + e + f = 0\). 记 \(r_1(A)\) 为 \(A\) 的第4行各数之和 \((i = 1,2)\)，\(e_j(A)\) 为 \(A\) 的第 \(j\) 列各数之和 \((j = 1,2,3)\)；记 \(k(A)\) 为 \([r_1(A)]\)，\([r_2(A)]\)，\([e_3(A)]\)，\([e_2(A)]\)，\([e_3(A)]\) 中的最小值. (1) 对如下数表 A，求 \( k(A) \) 的值: 1 1 -0.8 0.1 -0.3 -1 (2) 设数表 \(A\) 形如 1 1 -1 -2d d d -1 其中 \(-1 \leqslant d \leqslant 0\). 求 \(k(A)\) 的最大值; (3) 对所有满足性质 \(P\) 的 2 行 3 列的数表 \(A\)，求 \(k(A)\) 的最大值.
\end{problem}

